\documentclass[11pt,a4paper]{article}
\usepackage[utf8]{inputenc}
\usepackage[T1]{fontenc}
\usepackage[spanish]{babel}
\usepackage{amsmath,amssymb,mathtools}
\usepackage{geometry}
\usepackage{enumitem}
\usepackage{lmodern}
\usepackage{microtype}
\geometry{margin=2.4cm}
\setlength{\parindent}{0pt}
\setlength{\parskip}{0.45em}
\newcommand{\R}{\mathbb{R}}
\newcommand{\Ker}{\operatorname{Ker}}
\newcommand{\Ima}{\operatorname{Im}}
\newcommand{\rg}{\operatorname{rg}}
\newcommand{\tr}{\operatorname{tr}}
\newcommand{\Span}{\operatorname{L}}
\newcommand{\proy}{\operatorname{proy}}
\newcommand{\diag}{\operatorname{diag}}
\newcommand{\sg}{\operatorname{sg}}
\begin{document}

\begin{center}
{\LARGE\bfseries Álgebra Lineal}\\[0.4em]
{\Large Tema 6: El Espacio Vectorial Euclídeo --- Soluciones}
\end{center}

\section{Ejercicio 1}
\subsection{Solución}

La matriz de Gram es:

\[
\boxed{
G=
\begin{pmatrix}
1&0&1\\
0&2&1\\
1&1&2
\end{pmatrix}.
}
\]

Para $v=(1,1,1)$:

\[
\|v\|^2
=
v^tGv
=
9.
\]

Por tanto:

\[
\boxed{\|v\|=3.}
\]

Sean:

\[
u=(1,0,1),
\qquad
w=(0,1,0).
\]

\[
\langle u,w\rangle=1,
\qquad
\|u\|=\sqrt5,
\qquad
\|w\|=\sqrt2.
\]

Así:

\[
\boxed{
\cos\alpha=\frac1{\sqrt{10}},
\qquad
\alpha=\arccos\left(\frac1{\sqrt{10}}\right).
}
\]

Para $(x,y,z)\in S^\perp$:

\[
x+z=0,
\qquad
2y+z=0.
\]

Por tanto:

\[
\boxed{
S^\perp
=
\Span\{(-2,-1,2)\}.
}
\]

\section{Ejercicio 2}
\subsection{Solución}

Sea:

\[
p(x)=a+bx+cx^2.
\]

Si $p\in S\cap T$:

\[
a=0,
\qquad
b+c=0.
\]

Luego:

\[
\boxed{
S\cap T
=
\Span\{x-x^2\}.
}
\]

Por Grassmann:

\[
\dim(S+T)
=
\dim S+\dim T-\dim(S\cap T)
=
2+2-1=3.
\]

Así:

\[
\boxed{
S+T=\R_2[x].
}
\]

Tomando el suplementario ortogonal en la base canónica:

\[
(0,1,-1)\cdot(a,b,c)=0,
\]

por lo que:

\[
b-c=0.
\]

Así:

\[
\boxed{
(S\cap T)^\perp
=
\Span\{1,x+x^2\}
}
\]

con ecuación implícita:

\[
\boxed{b=c.}
\]

\section{Ejercicio 3}
\subsection{Solución}

Las ecuaciones de $W$ tienen vectores normales:

\[
n_1=(1,1,-1,1),
\qquad
n_2=(2,1,-1,3).
\]

Con el producto escalar usual:

\[
W^\perp
=
\Span\{n_1,n_2\}.
\]

Por tanto:

\[
\boxed{
W^\perp
=
\Span\{(1,1,-1,1),(2,1,-1,3)\}.
}
\]

\section{Ejercicio 4}
\subsection{Solución}

\[
\|u+v\|^2
=
\langle u+v,u+v\rangle
\]

\[
=
\|u\|^2+\|v\|^2+2\langle u,v\rangle.
\]

Por tanto:

\[
\|u+v\|^2
=
\|u\|^2+\|v\|^2
\]

si y sólo si:

\[
\boxed{\langle u,v\rangle=0.}
\]

\section{Ejercicio 5}
\subsection{Solución}

\[
\|u+v\|^2
=
\|u\|^2+\|v\|^2+2\langle u,v\rangle
\]

\[
\|u-v\|^2
=
\|u\|^2+\|v\|^2-2\langle u,v\rangle.
\]

Sumando:

\[
\boxed{
\|u+v\|^2+\|u-v\|^2
=
2\left(\|u\|^2+\|v\|^2\right).
}
\]

\section{Ejercicio 6}
\subsection{Solución}

\[
\langle u+v,u-v\rangle
=
\|u\|^2-\|v\|^2.
\]

Por tanto:

\[
u+v\perp u-v
\iff
\|u\|^2=\|v\|^2
\iff
\boxed{\|u\|=\|v\|.}
\]

En un paralelogramo, $u$ y $v$ representan dos lados y:

\[
u+v,
\qquad
u-v
\]

sus diagonales.

Por tanto:

\[
\boxed{
\text{es un rombo}
\iff
\text{sus diagonales son perpendiculares}.
}
\]

\section{Ejercicio 7}
\subsection{Solución}

Una base de $F$ es:

\[
\{(-1,1,0),(-2,0,1)\}.
\]

Sea:

\[
v=(x_1,x_2,x_3)\in F^\perp.
\]

Entonces:

\[
-x_1+2x_2=0
\]

\[
-2x_1+3x_3=0.
\]

Por tanto:

\[
\boxed{
F^\perp
=
\Span\{(6,3,4)\}.
}
\]

\section{Ejercicio 8}
\subsection{Solución}

Para el apartado a), imponemos:

\[
x+2y+z=0
\]

\[
-y+z=0
\]

\[
x+y-2z+t=0.
\]

Se obtiene:

\[
(x,y,z,t)
=
\lambda(-3,1,1,4).
\]

Como:

\[
\|(-3,1,1,4)\|^2=27,
\]

un vector unitario es:

\[
\boxed{
\frac1{\sqrt{27}}(-3,1,1,4).
}
\]

Para el apartado b), tomamos:

\[
v_1=(1,2,-1,0),
\quad
v_2=(1,0,-2,1),
\quad
v_3=(0,1,1,0).
\]

Gram--Schmidt da:

\[
u_1=(1,2,-1,0)
\]

\[
u_2=\frac12(1,-2,-3,2)
\]

\[
u_3=\frac19(1,1,3,5).
\]

Normalizando:

\[
\boxed{
e_1=
\frac1{\sqrt6}(1,2,-1,0)
}
\]

\[
\boxed{
e_2=
\frac1{\sqrt{18}}(1,-2,-3,2)
}
\]

\[
\boxed{
e_3=
\frac16(1,1,3,5).
}
\]

\section{Ejercicio 9}
\subsection{Solución}

Sea:

\[
s=
\alpha(0,1,2)
+
\beta(1,2,3).
\]

Exigimos:

\[
(1,2,1)-s
\perp
(0,1,2)
\]

y:

\[
(1,2,1)-s
\perp
(1,2,3).
\]

Se obtiene:

\[
5\alpha+8\beta=4
\]

\[
8\alpha+14\beta=8.
\]

Por tanto:

\[
\alpha=-\frac43,
\qquad
\beta=\frac43.
\]

Así:

\[
\boxed{
\proy_S(1,2,1)
=
\left(
\frac43,\frac43,\frac43
\right).
}
\]

\section{Ejercicio 10}
\subsection{Solución}

En la base $\{1,x\}$:

\[
\langle1,1\rangle=1,
\qquad
\langle1,x\rangle=\frac12,
\qquad
\langle x,x\rangle=\frac13.
\]

Por tanto:

\[
\boxed{
G=
\begin{pmatrix}
1&\frac12\\
\frac12&\frac13
\end{pmatrix}.
}
\]

Además:

\[
\langle x+3,2x+4\rangle
=
\frac{53}{3}
\]

\[
\|x+3\|^2=\frac{37}{3}
\]

\[
\|2x+4\|^2=\frac{76}{3}.
\]

Así:

\[
\boxed{
\cos\alpha
=
\frac{53}{\sqrt{2812}}.
}
\]

Para la proyección sobre $x+2$:

\[
\frac{\langle x+3,x+2\rangle}
{\|x+2\|^2}
=
\frac{53}{38}.
\]

Por tanto:

\[
\boxed{
\proy_{x+2}(x+3)
=
\frac{53}{38}(x+2).
}
\]

Aplicando Gram--Schmidt a $\{1,x\}$:

\[
u_1=1
\]

\[
u_2=x-\frac12.
\]

Como:

\[
\left\|x-\frac12\right\|^2
=
\frac1{12},
\]

una base ortonormal es:

\[
\boxed{
\left\{
1,\;
2\sqrt3\,x-\sqrt3
\right\}.
}
\]

\section{Ejercicio 11}
\subsection{Solución}

Como:

\[
S=
\Span\{(1,0,0),(0,1,0)\},
\]

se tiene:

\[
\boxed{
\proy_S(1,1,1)
=
(1,1,0).
}
\]

\section{Ejercicio 12}
\subsection{Solución}

Para el ajuste cuadrático:

\[
y=a_0+a_1x+a_2x^2.
\]

Las ecuaciones normales dan:

\[
a_0=-\frac{11}{4},
\qquad
a_1=\frac{79}{20},
\qquad
a_2=-\frac14.
\]

Por tanto:

\[
\boxed{
y=
-\frac{11}{4}
+
\frac{79}{20}x
-
\frac14x^2.
}
\]

Para el ajuste cúbico, al haber cuatro puntos se obtiene el polinomio interpolador:

\[
\boxed{
y=
-1
+
\frac76x
+
x^2
-
\frac16x^3.
}
\]

\section{Ejercicio 13}
\subsection{Solución}

Buscamos:

\[
y=a+bx.
\]

Las ecuaciones normales son:

\[
\begin{pmatrix}
4&\frac{43}{10}\\
\frac{43}{10}&\frac{477}{100}
\end{pmatrix}
\begin{pmatrix}
a\\b
\end{pmatrix}
=
\begin{pmatrix}
11\\
\frac{127}{10}
\end{pmatrix}.
\]

Se obtiene:

\[
a=-\frac{300}{59},
\qquad
b=\frac{430}{59}.
\]

Por tanto:

\[
\boxed{
y=
-\frac{300}{59}
+
\frac{430}{59}x.
}
\]

\section{Ejercicio 14}
\subsection{Solución}

El ajuste lineal de:

\[
(1,4),(2,8),(3,11),(4,14)
\]

es:

\[
\boxed{
y=1+\frac{33}{10}x.
}
\]

Para:

\[
x=7,
\]

\[
y(7)
=
1+\frac{231}{10}
=
\frac{241}{10}.
\]

Por tanto:

\[
\boxed{
N(7)\approx24.1
\text{ miles de microorganismos}.
}
\]

\section{Ejercicio 15}
\subsection{Solución}

El polinomio característico es:

\[
p(\lambda)
=
\lambda(1-\lambda)(\lambda-3).
\]

Los autovalores son:

\[
0,\quad1,\quad3.
\]

Podemos tomar:

\[
u_1=(-1,1,-1),
\quad
u_2=(1,0,-1),
\quad
u_3=(1,2,1).
\]

Son ortogonales.

Normalizando:

\[
e_1=
\frac1{\sqrt3}(-1,1,-1)
\]

\[
e_2=
\frac1{\sqrt2}(1,0,-1)
\]

\[
e_3=
\frac1{\sqrt6}(1,2,1).
\]

Por tanto:

\[
\boxed{
P=
\begin{pmatrix}
-\frac1{\sqrt3}&\frac1{\sqrt2}&\frac1{\sqrt6}\\
\frac1{\sqrt3}&0&\frac2{\sqrt6}\\
-\frac1{\sqrt3}&-\frac1{\sqrt2}&\frac1{\sqrt6}
\end{pmatrix}
}
\]

y:

\[
\boxed{
P^tAP=
\diag(0,1,3).
}
\]

\section{Ejercicio 16}
\subsection{Solución}

Las coordenadas de:

\[
2e_1+e_2
\]

en $B$ son:

\[
(1,1,0).
\]

Entonces:

\[
\langle2e_1+e_2,2e_1+e_2\rangle
=
(1,1,0)
G
\begin{pmatrix}
1\\1\\0
\end{pmatrix}
=3.
\]

Así:

\[
\boxed{\text{1) Falsa}.}
\]

Además:

\[
[e_1]_B=(1,0,0),
\]

\[
[e_1+e_2]_B=(0,1,0),
\]

\[
[e_3]_B=(0,-1,1).
\]

Se verifica:

\[
\langle e_1,e_1+e_2\rangle=0
\]

y:

\[
\langle e_1,e_3\rangle=0.
\]

Por tanto:

\[
\boxed{\text{2) Cierta}.}
\]

La matriz de cambio es:

\[
P=
\begin{pmatrix}
1&-1&0\\
0&1&-1\\
0&0&1
\end{pmatrix}.
\]

Así:

\[
G_{B'}
=
P^tGP
=
\begin{pmatrix}
1&-1&0\\
-1&3&-2\\
0&-2&3
\end{pmatrix}.
\]

Por tanto:

\[
\boxed{\text{3) Falsa}.}
\]

Si $Q$ es ortogonal:

\[
Q^tQ=I.
\]

Tomando determinantes:

\[
(\det Q)^2=1.
\]

Por tanto:

\[
\boxed{\text{4) Cierta}.}
\]

La afirmación 5 no es cierta en general para matrices ortogonales reales: sus autovalores complejos pueden tener módulo $1$ sin ser $\pm1$.

Por tanto:

\[
\boxed{\text{5) Falsa}.}
\]

\section{Ejercicio 17}
\subsection{Solución}

Sean $A,B$ ortogonales:

\[
A^tA=I,
\qquad
B^tB=I.
\]

Entonces:

\[
(AB)^tAB
=
B^tA^tAB
=
B^tB
=
I.
\]

Por tanto:

\[
\boxed{AB\text{ es ortogonal}.}
\]

Además:

\[
A^{-1}=A^t.
\]

Entonces:

\[
(A^{-1})^tA^{-1}
=
AA^t
=
I.
\]

Por tanto:

\[
\boxed{A^{-1}\text{ es ortogonal}.}
\]

\section{Ejercicio 18}
\subsection{Solución}

Las columnas deben ser unitarias.

Para la primera:

\[
\frac{4+9+s^2}{49}=1
\]

de donde:

\[
s=\pm6.
\]

Para la segunda:

\[
\frac{t^2+s^2+4}{49}=1.
\]

Así:

\[
t=\pm3.
\]

Imponiendo además la ortogonalidad entre columnas se obtiene:

\[
\boxed{
s=6,
\qquad
t=-3.
}
\]

\section{Ejercicio 19}
\subsection{Solución}

Sean:

\[
C_1,\ldots,C_n
\]

las columnas de $A$.

Entonces:

\[
A^tA
=
\left(
C_i^tC_j
\right)
=
\left(
\langle C_i,C_j\rangle
\right).
\]

Así:

\[
A^tA=I
\]

si y sólo si:

\[
\langle C_i,C_j\rangle
=
\begin{cases}
1,&i=j,\\
0,&i\neq j.
\end{cases}
\]

Por tanto:

\[
\boxed{
A\text{ es ortogonal}
\iff
\{C_1,\ldots,C_n\}
\text{ es ortonormal}.
}
\]

\section{Ejercicio 20}
\subsection{Solución}

El polinomio característico es:

\[
p(\lambda)
=
-(1+\lambda)^2(\lambda-3).
\]

Los autovalores son:

\[
-1\text{ doble},
\qquad
3\text{ simple}.
\]

Para $\lambda=-1$ podemos tomar:

\[
u_1=(1,0,-1),
\qquad
u_2=(0,1,0).
\]

Para $\lambda=3$:

\[
u_3=(1,0,1).
\]

Normalizando:

\[
\boxed{
P=
\begin{pmatrix}
\frac1{\sqrt2}&0&\frac1{\sqrt2}\\
0&1&0\\
-\frac1{\sqrt2}&0&\frac1{\sqrt2}
\end{pmatrix}
}
\]

y:

\[
\boxed{
P^tAP
=
\diag(-1,-1,3).
}
\]

\section{Ejercicio 21}
\subsection{Solución}

Para $A$:

\[
p_A(\lambda)
=
-\lambda^3+9\lambda^2-16\lambda
\]

tiene un autovalor nulo y dos positivos.

Por tanto:

\[
\sg(A)=(2,0).
\]

Para $B$:

\[
p_B(\lambda)
=
-\lambda^3+10\lambda^2+8\lambda
\]

tiene un autovalor nulo, uno positivo y uno negativo.

Así:

\[
\sg(B)=(1,1).
\]

Como dos matrices simétricas reales son congruentes si y sólo si tienen la misma inercia:

\[
\boxed{
A\text{ y }B\text{ no son congruentes}.
}
\]

\section{Ejercicio 22}
\subsection{Solución}

Los autovalores son:

\[
3,
\qquad
-1.
\]

Podemos tomar:

\[
v_1=(-1,1),
\qquad
v_2=(1,1).
\]

Normalizando:

\[
e_1=
\frac1{\sqrt2}(-1,1),
\qquad
e_2=
\frac1{\sqrt2}(1,1).
\]

Por tanto:

\[
\boxed{
P=
\frac1{\sqrt2}
\begin{pmatrix}
-1&1\\
1&1
\end{pmatrix}
}
\]

y:

\[
\boxed{
P^tMP=
\diag(3,-1).
}
\]

\section{Ejercicio 23}
\subsection{Solución}

Los autovalores son:

\[
-1\text{ doble},
\qquad
2\text{ simple}.
\]

Una base ortogonal de autovectores es:

\[
u_1=(1,0,-1)
\]

\[
u_2=(1,-2,1)
\]

\[
u_3=(1,1,1).
\]

Normalizando:

\[
\boxed{
P=
\begin{pmatrix}
\frac1{\sqrt2}&\frac1{\sqrt6}&\frac1{\sqrt3}\\
0&-\frac2{\sqrt6}&\frac1{\sqrt3}\\
-\frac1{\sqrt2}&\frac1{\sqrt6}&\frac1{\sqrt3}
\end{pmatrix}
}
\]

y:

\[
\boxed{
P^tAP
=
\diag(-1,-1,2).
}
\]

\section{Ejercicio 24}
\subsection{Solución}

El plano puede escribirse como:

\[
S=
\{(-2\alpha,\beta,\alpha):
\alpha,\beta\in\R\}.
\]

Sea:

\[
s=(-2\alpha,\beta,\alpha)
\]

la proyección de:

\[
v=(1,1,1).
\]

Entonces:

\[
v-s
=
(1+2\alpha,1-\beta,1-\alpha)
\]

debe ser ortogonal a:

\[
(-2,0,1)
\]

y:

\[
(0,1,0)
\]

respecto del producto de Gram $G$.

Se obtiene:

\[
1-20\alpha-5\beta=0
\]

\[
1-5\alpha-2\beta=0.
\]

Resolviendo:

\[
\alpha=-\frac15,
\qquad
\beta=1.
\]

Por tanto:

\[
\boxed{
\proy_S(1,1,1)
=
\left(
\frac25,1,-\frac15
\right).
}
\]

\section{Ejercicio 25}
\subsection{Solución}

Como $A$ es real y simétrica, es ortogonalmente diagonalizable.

Sus autovalores son:

\[
1\text{ doble},
\qquad
4\text{ simple}.
\]

Una base ortogonal de autovectores es:

\[
u_1=(1,0,-1),
\qquad
u_2=(1,-2,1),
\qquad
u_3=(1,1,1).
\]

Se verifica:

\[
u_1\cdot u_2
=
u_1\cdot u_3
=
u_2\cdot u_3
=
0.
\]

Normalizando:

\[
\boxed{
B'=
\left\{
\frac1{\sqrt2}(1,0,-1),
\frac1{\sqrt6}(1,-2,1),
\frac1{\sqrt3}(1,1,1)
\right\}.
}
\]

La matriz de paso es:

\[
\boxed{
P=
\begin{pmatrix}
\frac1{\sqrt2}&\frac1{\sqrt6}&\frac1{\sqrt3}\\
0&-\frac2{\sqrt6}&\frac1{\sqrt3}\\
-\frac1{\sqrt2}&\frac1{\sqrt6}&\frac1{\sqrt3}
\end{pmatrix}
}
\]

y:

\[
P^tP=I.
\]

Por tanto:

\[
\boxed{P\text{ es ortogonal}.}
\]

Además:

\[
\boxed{
P^tAP=
\diag(1,1,4).
}
\]

Todos los autovalores son positivos, luego:

\[
\boxed{
q(x)=x^tAx
\text{ es definida positiva}.
}
\]

\section{Ejercicio 26}
\subsection{Solución}

En:

\[
B=\{(1,1),(1,2)\}
\]

se tiene:

\[
f(1,1)
=
-3(1,1)+2(1,2)
\]

y:

\[
f(1,2)
=
-4(1,1)+3(1,2).
\]

Por tanto:

\[
A=
\begin{pmatrix}
-3&-4\\
2&3
\end{pmatrix}.
\]

La matriz de Gram del producto usual en $B$ es:

\[
G_B=
\begin{pmatrix}
2&3\\
3&5
\end{pmatrix}.
\]

Calculando:

\[
A^tG_BA
=
\begin{pmatrix}
2&3\\
3&5
\end{pmatrix}.
\]

Así:

\[
\boxed{
G_B=A^tG_BA.
}
\]

\section{Ejercicio 27}
\subsection{Solución}

Como:

\[
S\subseteq S(a)
\]

y:

\[
\dim S=2,
\]

el autovalor $a$ tiene multiplicidad al menos $2$.

Si:

\[
a\neq b,
\]

los autovalores son:

\[
a,a,b.
\]

Una base ortonormal adaptada es:

\[
e_1=
\frac1{\sqrt2}(1,0,-1)
\]

\[
e_2=
\frac1{\sqrt6}(1,-2,1)
\]

\[
e_3=
\frac1{\sqrt3}(1,1,1).
\]

Así:

\[
\boxed{
P=
\begin{pmatrix}
\frac1{\sqrt2}&\frac1{\sqrt6}&\frac1{\sqrt3}\\
0&-\frac2{\sqrt6}&\frac1{\sqrt3}\\
-\frac1{\sqrt2}&\frac1{\sqrt6}&\frac1{\sqrt3}
\end{pmatrix}
}
\]

y:

\[
\boxed{
J=
\diag(a,a,b).
}
\]

Si:

\[
a=b,
\]

todos los autovalores son $a$ y:

\[
\boxed{
A=aI,
\qquad
J=aI,
\qquad
P=I.
}
\]

Para que $A$ sea matriz de un producto escalar debe ser definida positiva.

Por tanto:

\[
\boxed{
a>0,
\qquad
b>0.
}
\]

Los valores naturales más pequeños son:

\[
a=b=1.
\]

Entonces:

\[
A=I.
\]

Para:

\[
u=(1,-1,0),
\qquad
v=(1,1,1),
\]

\[
\langle u,v\rangle=0.
\]

Por tanto:

\[
\boxed{
\alpha=\frac{\pi}{2}.
}
\]

\section{Ejercicio 28}
\subsection{Solución}

Para:

\[
u=(1,1,1),
\]

\[
U^tU=3.
\]

Por tanto:

\[
H
=
I-\frac23UU^t.
\]

Así:

\[
\boxed{
H=
\frac13
\begin{pmatrix}
1&-2&-2\\
-2&1&-2\\
-2&-2&1
\end{pmatrix}.
}
\]

Aplicando $H$:

\[
\boxed{
H(3,3,3)=(-3,-3,-3)
}
\]

\[
\boxed{
H(1,-1,0)=(1,-1,0)
}
\]

\[
\boxed{
H(1,2,3)=(-3,-2,-1).
}
\]

\section{Ejercicio 29}
\subsection{Solución}

Como:

\[
B'=\{e_1,e_2+e_3,e_2-e_3\}
\]

es ortogonal:

\[
\langle e_1,e_2\rangle
=
\langle e_1,e_3\rangle
=
0
\]

y:

\[
\|e_2\|=\|e_3\|.
\]

Además:

\[
\cos\frac{\pi}{3}
=
\frac{\langle e_2,e_3\rangle}
{\|e_2\|\|e_3\|}
=
\frac12.
\]

Los números naturales pares más pequeños son:

\[
\|e_1\|=\|e_2\|=\|e_3\|=2.
\]

Así:

\[
\langle e_2,e_3\rangle=2.
\]

Por tanto:

\[
\boxed{
G=
\begin{pmatrix}
4&0&0\\
0&4&2\\
0&2&4
\end{pmatrix}.
}
\]

Para:

\[
S=
\Span\{e_1+e_2+e_3,e_1+2e_2\},
\]

sea:

\[
u=xe_1+ye_2+ze_3\in S^\perp.
\]

Las condiciones son:

\[
4x+6y+6z=0
\]

\[
4x+8y+4z=0.
\]

Se obtiene:

\[
(x,y,z)=\lambda(-3,1,1).
\]

Como:

\[
\|(-3,1,1)\|^2=48,
\]

una base ortonormal es:

\[
\boxed{
\left\{
\frac1{\sqrt{48}}(-3e_1+e_2+e_3)
\right\}.
}
\]

Además:

\[
e_2-e_3\in S.
\]

Por tanto:

\[
\boxed{
\proy_{S^\perp}(e_2-e_3)=0.
}
\]

\section{Ejercicio 30}
\subsection{Solución}

La matriz asociada es:

\[
A=
\begin{pmatrix}
10&-6\\
-6&5
\end{pmatrix}.
\]

El polinomio característico es:

\[
\lambda^2-15\lambda+14.
\]

Por tanto:

\[
\lambda_1=14,
\qquad
\lambda_2=1.
\]

Autovectores asociados:

\[
v_1=(3,-2),
\qquad
v_2=(2,3).
\]

Normalizando:

\[
e_1=
\frac1{\sqrt{13}}(3,-2),
\qquad
e_2=
\frac1{\sqrt{13}}(2,3).
\]

Así:

\[
\boxed{
B'=
\left\{
\frac1{\sqrt{13}}(3,-2),
\frac1{\sqrt{13}}(2,3)
\right\}
}
\]

y:

\[
\boxed{
P=
\frac1{\sqrt{13}}
\begin{pmatrix}
3&2\\
-2&3
\end{pmatrix}.
}
\]

Por tanto:

\[
\boxed{
P^tAP=
\diag(14,1).
}
\]

En la nueva base:

\[
\boxed{
14x'^2+y'^2=1.
}
\]

\section{Ejercicio 31}
\subsection{Solución}

Las coordenadas de:

\[
2e_1+e_2
\]

en $B$ son:

\[
(1,1,0).
\]

Por tanto:

\[
\boxed{
\langle2e_1+e_2,2e_1+e_2\rangle
=
3.
}
\]

Las coordenadas en $B$ de los vectores de $B'$ son:

\[
[e_1]_B=(1,0,0)
\]

\[
[e_2]_B=(-1,1,0)
\]

\[
[e_3]_B=(0,-1,1).
\]

Luego:

\[
P=
\begin{pmatrix}
1&-1&0\\
0&1&-1\\
0&0&1
\end{pmatrix}.
\]

Así:

\[
G_{B'}
=
P^tGP.
\]

Por tanto:

\[
\boxed{
G_{B'}
=
\begin{pmatrix}
1&-1&0\\
-1&3&-2\\
0&-2&3
\end{pmatrix}.
}
\]

\[
\langle e_1,e_1+e_2\rangle=0
\]

y:

\[
\langle e_1,e_3\rangle=0.
\]

Por tanto:

\[
\boxed{
V_1\perp V_2.
}
\]

\end{document}
